\section{Organisation and Module Walkthrough}

\begin{frame}{Introduction}
    \mycomment{I am honoured to teach you this semester.  I have an extensive background in academia and as a consultant and manager.}
    \mytext{
    Jan Suchanek
    \begin{itemize}
        \item Diplom-Volkswirt (MA Economics)
        \item MBL (Master of Business Law)
        \item based in Berlin, born in Hamburg
        \item English, German (native), French, Spanish
        \item work as consultant, project manager, lecturer
        \item Project Manager: Energy Networking Projects
        \item Teaching Project Management
        \item Teaching Foreign Trade (ISGM)
        \item ...and this class
    \end{itemize}
    }
\end{frame}

\begin{frame}{Introduction}
    \mytext{
    Attendants:
        \begin{itemize}
            \item BIM students
            \item Incoming Students?
            \item favourite market?
        \end{itemize}
    }
\end{frame}

\begin{frame}{Schedule, Grading and Procedure}
    \mycomment{This class is predominantly an online class. There are, however, some sessions held in situ.

    Lectures will cover the theory. Student presentations shall provide the illustration from real life applications by outlining the most important properties of specific markets. 
    
    }
    \mytext{
        \begin{itemize}
            \item April -- July: online lectures
            \item Mai 18th,  Room D207
            \item[] Student Present a Market (early bird)
            \item June 26th \& June 27th (Saturday!) Room D107
            \item[] Student Present a Market (last chance)
            \item End July: Written Exam
        \end{itemize}

        % \medskip
        
        \begin{description}
            \item[50\%] Student Presentation 
            \item[50\%] digitally written exam 
        \end{description}
    }
\end{frame}

\begin{frame}{Schedule, Grading and Procedure}
    \mytext{
        \begin{itemize}
            \item Students present a market
            \item groups per 1-5 students
            \item 15 minutes presentation, 5 minutes discussion
            \item Two options:
            \begin{enumerate}
                \item Early bird -- Mai 18th (max. 8 gr.)
                \item Last chance -- June 26th and 27th
            \end{enumerate}
            \item state your intentions on forum (\href{}{LINK})
        \end{itemize}

    }
    
    \mycomment{
        Every student must speak over one slide / cover one aspect. 
        
        Grading will be 50\% overall quality and consistence of the presentation and 50\% individual performance.    
    } 
    
\end{frame}



\begin{frame}{Economic Theory}
    \showmytikz{../shared/tikz/tikz-tree1-micromacro.tex}
    \mycomment{
        \begin{itemize}
            \item Micro-Economics looks at individuals
            \item Micro-Economics looks at single transactions
            \item Micro-Economics looks at single markets
            \item Micro Economics looks at \textit{partial equilibria}, i.e. equilibria of single markets
            \item Micro-Economics are foundational for Economics
            \item Macro-Economics looks at the economy: aggregates
            \item Macro-Economics looks at the \textit{total equilibrium}, e.g. of the economy as a whole.
        \end{itemize}
        }
\end{frame}

\begin{frame}{Microeconomics within Economics}
    \showmytikz{../shared/tikz/tikz-tree2-micromacro.tex}
        \mycomment{
            Micro-economics looks at the fundamentals of economic behaviour under idealised circumstances.

            
            Micro-economics looks at isolated (and idealised) individual effects, not at the economy as an entire system.            
        }
\end{frame}
\begin{frame}{Macroeconomics within Economics}
    \showmytikz{../shared/tikz/tikz-tree3-micromacro.tex}
        \mycomment{
            Macroeconomics looks at the economy as a circular system. There are topics belonging to both branches, and Microeconomics provide the fundamentals for Macroeconomics. However, the sum is more than its parts.         
        }
\end{frame}

\begin{frame}{Economic Theory}
    \mycomment{
        In physics, the falling of objects can be best understood by abstracting from the presence of air resistance to isolate the pure effect of gravity. 
        
        Economics takes a similar approach. Starting from first principles, Microeconomics analyses optimal choices and their welfare results in a generalised manner.
                
        % There is also the empiric branch, and applied Economics.
        
        Among empirical economics is econometrics, statistics, also modelling (of concrete cases), hence forecasting, among with advice and consultancy to public and private entities, and the study of economic history.

        Of course, theory is helpful for empirical studies, and for policy analysis and normative recommendations.
    }
    \mytext{
        \begin{center}
            Micro-Economics is in Economics 
            
            \textit{almost like}
            
            Theoretical Physics in Physics.
        \end{center}

        
        \begin{itemize} 
            \item fundamental descriptions 
            \item high level of abstraction
            \item high level of simplification
            \item isolation of effects
            \item fundamental insights
            \item basis for more sophistication
        \end{itemize}
    }
\end{frame}

\begin{frame}{Aspects we leave to Macroeconomics}
    \mycomment{
        Everything that concerns \textit{the economy as such} or \textit{as a whole} is the field of Macroeconomics.
    }
    \mytext{
        \begin{itemize}
            \item circuits, entire economies, gross domestic product, savings and investment
            \item public finance, economic policy as aggregates
            \item taxes, subsidies, development
            \item money, credit, currencies, stocks
            \item innovation, intellectual property, education, growth
            \item business cycles, bubbles, (financial) crises
            \item resources, environment
            \item employment, equality, power and ownership
            \item trade, logistics, international balances, tariffs
        \end{itemize}
        }
\end{frame}

\begin{frame}{Circular Model of the Economy}
    \showmytikz[1,2]{../shared/tikz/tikz-flowmodel-own}
    \mycomment{
        We distinguish two roles: households, and businesses. Households sell work and labour to businesses against wages and salaries on the labour markets, so that they be able to buy goods and services in return.

        So while consumers (households) make consumption decisions, businesses decide what to produce.

    }
\end{frame}
\begin{frame}{Circular Model of the Economy}
    \showmytikz[1,2,3]{../shared/tikz/tikz-flowmodel-own}
    \mycomment{
        Households save (on capital markets), their savings become investments (of business). However, these capital markets are more a concern of Macroeconomics.
    }
\end{frame}
\begin{frame}{Circular Model of the Economy}
    \showmytikz[1,2,3,4]{../shared/tikz/tikz-flowmodel-own}
    \mycomment{
        On top, households pay taxes (to the government), which the government spends on markets (to businesses). 

        The government also acts as rule giver and regulator on markets, and can impose levies, taxes, and tariffs on markets as well, such as value added tax (VAT).
        
        Households also import goods and services from the rest of the world, while business export to the rest of the world.

        % Note that this is simplified. The government may also redistribute, and has other functions as well. Households decide how much time (and skill) to offer on labour markets. There is also an international capital market. Businesses may be organised in one way or the other, and households may be large or small.
    }
\end{frame}

\begin{frame}{Microeconomics within Economics}
    \showmytikz{../shared/tikz/tikz-tree2-micromacro-markets.tex}
        \mycomment{
            Markets take centre stage. We are looking at everything else in respect to markets, really.

            Our main tools will be mathematical models. The next section shall give you a preview---we will get into the sense of all this much deeper later.
        }
\end{frame}

\begin{frame}{Recap: Mathematical Optimisation}
    \mycomment{
        Micro-Economics looks at individual markets and at actors making choices on those markets. 
        
        % We assume e to maximise their utility; and that they do this by maximising their income and minimising their expense.

        For example, I can offer to teach classes for a particular price. If I set my price very high (P), nobody will book me. If I do it for free, I will be booked (S) times.

        So how much I will be booked (volume $x$) will be a decreasing function of my price $p$.



        \begin{equation}
            \frac{d p}{d x} <0
        \end{equation}

        Here: 
        \begin{equation}
            f(x) = p(x) = -ax + b
        \end{equation}

    }
    \showmytikz[]{../shared/tikz/tikz-optimisation.tex}
\end{frame}
\begin{frame}{Recap: Mathematical Optimisation}
    \mycomment{
        Now assume that I am trying to earn as much money as I can. Which price should I choose?

        I want to optimise my turnover which here, assuming that I don't have any costs, is my profit.

        And this turns out to be volume times price.
        
        \begin{equation}
            g(x) = p(x)\cdot x            
        \end{equation}

    }
    \showmytikz[1]{../shared/tikz/tikz-optimisation.tex}
\end{frame}
\begin{frame}{Recap: Mathematical Optimisation}
    \mycomment{
        So here we ask: which combination of $p$ and $x$ will give us the largest rectangle under the curve? This rectangle is my profit (turnover).

    }
    \showmytikz[1,2]{../shared/tikz/tikz-optimisation.tex}
\end{frame}
\begin{frame}{Recap: Mathematical Optimisation}
    \mycomment{
        The size of the rectangle will be at that combination of $x,p$ where 

        \begin{equation}
            g(x) = p(x)\cdot x \rightarrow max!   
        \end{equation}
        
    }
    \showmytikz[1,2,3]{../shared/tikz/tikz-optimisation.tex}
\end{frame}
\begin{frame}{Recap: Mathematical Optimisation}
    \mycomment{
        The function we need to optimise is a parabola, because the rectangle surface is a parabola (assuming a linear form of $p(x)$).

        \begin{equation}
            g(x) = p(x)\cdot x            
        \end{equation}

        (Parabola is not drawn to scale.)
        
        You remember that the optimum can be found by setting the derivative zero and solving for $x$.

        \begin{equation}
            g'(x)\stackrel{!}{=}0 
        \end{equation}

    }
    \showmytikz[4]{../shared/tikz/tikz-optimisation.tex}
\end{frame}
\begin{frame}{Recap: Mathematical Optimisation}
    \mycomment{
        What is the condition for the profit maximising price in this simple model?
        
        To find the actual solution to this model is left to the student as an exercise.

        \begin{equation}
            p(x) = -ax+b
        \end{equation}
        \begin{equation}
            g(x) = \left(-ax+b\right)x \rightarrow max!
        \end{equation}
    }
    \showmytikz[3,4]{../shared/tikz/tikz-optimisation.tex}
\end{frame}
\begin{frame}{Recap: Mathematical Optimisation}
    \showmytikz[0]{../shared/tikz/tikz-optimisationunderconstraints.tex}
    \mycomment{
        Sometimes we may have to optimise more complex functions under constraints.

    }
\end{frame}
\begin{frame}{Recap: Mathematical Optimisation}
    \showmytikz[1]{../shared/tikz/tikz-optimisationunderconstraints.tex}
    \mycomment{
        When constraints exist, it is important to check if these constraints are maxima. That means that the actor in question cannot \textit{really} optimise. 

        For example, anyone with a very low income will not really be able to choose how to spend this income best (but spend it all on rice and rent).
        
        The optimisation of more complex functions under constraints requires Lagrange multipliers.
    }
\end{frame}

\begin{frame}{Recap: Mathematical Optimisation}
    \mytext{
        The models we will use in this class are deliberately extremely simple. We will demonstrate all results graphically and not require algebra.    
    }
\end{frame}

\begin{frame}{Questions we will address}
    \mytext{
    Depending on the market:
    \begin{itemize}
        \item What exactly is that function $p(x)$?
        \item What else influences the choices?
        \begin{itemize}
            \item Monopoly or competition?
            \item Costs of production?
        \end{itemize}
        \item Who gains: buyers, or sellers, or both?
        \item Who will produce what, and how much?
        \item What goods and services will people buy?
        \item How are things distributed?
        \item How can the state influence the outcomes?
    \end{itemize}
    }
    \mycomment{
    These questions - and some more - will be addressed as we dive deeper into economic modelling.
    }
\end{frame}

\begin{frame}{Students Presentation of the Market for a Product}
    \mycomment{
    You may treat some of these aspects jointly unless your group is large enough to cover these all.

    We will of course cover the meaning of these terms in the lecture.
    }
    \mytext{
    Present the market of a particular good:
    \begin{enumerate}
        \item delimitation of the market
        \item volumes, market places        
        \item elasticities, substitutes
        \item connected markets
        \item buyers, sellers, regulators
        \item trends (technological)
        \item global connections
        \item externalities    
    \end{enumerate}
    }
\end{frame}

\begin{frame}{Some Interesting Markets/Goods}
    \mycomment{
        As you research markets, you can draw on studies and publications. In most markets, companies, associations and organisations publish valuable data. 
        
        Trade fairs and their list of participants can also offer a good view into markets. 
        
        International organisations such as the EU commission, research institutes, universities and NGOs also publish valuable information.
    }
    \mytext{
    \mylonglist{
    \begin{itemize}
        \item any mineral
        \begin{itemize}
            \item crude oil
            \item helium
            \item lithium
            \item ...
        \end{itemize}
        \item particular (!) fast moving consumer good (FMCGs)
        \item gas turbines, electricity, water...
        % \item public transport in Bremen
        \item radio spectrums
        \item Bremen suburban housing
        \item particular skilled labour
        \item asset management software
        \item nutmeg
    \end{itemize}
    }
    }
\end{frame}

\begin{frame}{Some Interesting Markets/Goods}
    \mycomment{  
    }   
    \mytext{    
    \mylonglist{
        \begin{itemize}
            \item grains
            \begin{itemize}
                \item wheat
                \item rice
                \item soy
            \end{itemize}
            \item cotton
            \item gum arabic
            \item precision gear wheels
            \item household waste
            \item certificates of origin (CoOs)
            \item EU ETS \ce{CO2} emission rights
            \item Japanese government bonds  (JGBs)
            \item microchips
            \item shipping containers
            \item university graduates (labour)
        \end{itemize}
        }
    }
\end{frame}

\begin{frame}{Technical Requirements}
    \mycomment{
    Course language is English. Communication in any other language will be ignored. 

    Abstain from sending obscure proprietary or unnecessarily bloated formats. NEVER send attachments. Upload only .pdf.

    For email enquiries, use your HSB email only: requests from obscure email addresses will be treated as spam.

    The final exam will be held in electronic form.
    }
    \mytext{
    \begin{itemize}
        % \item students grouped by lecturer
        \item work in groups of 1-5 students
        \item every student presents exactly one aspect on one slide
        \item create consistency in design and content
        \item upload \textbf{as .pdf} on the day
        \item also bring .pdf on a USB stick
        \item grading is 50\% individual performance and 50\% group work quality.
    \end{itemize}
    }
\end{frame}

